\documentclass[11pt]{article}
\usepackage{mathblatt}
\begin{document}
\blattkopf*{Basisaufgaben}{BAS-Z8}{Basisaufgaben · Zettel · BAS-Z8}
\einheitenkopf[][Zettel 8]{Basisaufgaben · Zettel 8}
\zweigzeile{ohne Taschenrechner · je Aufgabe 1 Punkt · Lösungen auf der Rückseite}
% \small: zehn Aufgaben mit bis zu zwei Koordinatensystemen auf einer Seite (Render 28.09.)
\small

% 1: Bruchteil einer Fläche bestimmen – brueche-dezimalzahlen-basis-k1-v7 (Original 2026-FOR-B1b)
\begin{aufgabe}{Bei welcher Figur ist genau $\frac{1}{8}$ grau? \hfill (P26F) \\ \kreuz{P} \kreuz{Q} \kreuz{R} \kreuz{S}}

P \bruchrechteck[2.5]{1}{6} \quad Q \bruchkreis[0.8]{1}{4} \quad R \kreissektor[0.8]{45}{} \quad S \bruchrechteck[2.5]{1}{10}
\end{aufgabe}

% 2: Zahlen in verschiedenen Darstellungen vergleichen – brueche-dezimalzahlen-basis-k4-v7 (Original 2023-OS-B1f)
\begin{aufgabe}{Welcher Wert ist der größte: $0{,}9^2$; $0{,}85$; $88\,\%$; $\frac{7}{8}$? \hfill (P23) \\ \leerfeld}
\end{aufgabe}

% 3: Termwert berechnen – rationale-zahlen-basis-k1-v5 (Original 2026-FOR-B1g)
\begin{aufgabe}{Berechne den Wert von $-2 \cdot (x - 3)$ für $x = -1$. \hfill (P26F) \\ \leerfeld}
\end{aufgabe}

% 4: Vorzeichenregel anwenden – rationale-zahlen-basis-k2-v1 (Original 2015-OS-B1g)
\begin{aufgabe}{Eine negative Zahl wird mit einer positiven Zahl multipliziert. Das Ergebnis ist … \hfill (P15) \\ \kreuz{immer positiv} \kreuz{immer negativ} \kreuz{Das kann man nicht entscheiden.}}
\end{aufgabe}

% 5: Prozentwert berechnen – prozentrechnung-basis-k4-v5 (Original 2026-FOR-B1a)
\begin{aufgabe}{$40\,\%$ von $35\,€$ \hfill (P26F) \\ \leerfeld[€]}
\end{aufgabe}

% 6: Wert nach prozentualer Erhöhung berechnen – prozentrechnung-basis-k5-v1 (Original 2024-OS-B1e)
\begin{aufgabe}{Eine Kinokarte kostet $4\,€$ und wird um $25\,\%$ teurer. Wie hoch ist der neue Preis? \hfill (P24) \\ \leerfeld[€]}
\end{aufgabe}

% 7: Term zu Sachtext angeben – terme-basis-k2-v4 (Original 2023-OS-B1h)
\begin{aufgabe}{Das Sechsfache einer Zahl $b$ wird um 9 vergrößert. Welcher Term passt? \hfill (P23) \\ \kreuz{$6 \cdot (b + 9)$} \kreuz{$6b + 9$} \kreuz{$9b + 6$} \kreuz{$6 + 9b$}}
\end{aufgabe}

% 8: Winkel an geschnittenen Parallelen bestimmen – winkel-dreiecke-basis-k3-v2 (Original 2015-OS-B1d)
\begin{aufgabe}{\begin{minipage}[t]{0.6\linewidth}Die beiden waagerechten Geraden sind parallel. Wie groß ist $\alpha$? \hfill (P15) \\ $\alpha =$ \leerfeld °\end{minipage}\hfill\begin{minipage}[t]{0.37\linewidth}\vspace{0pt}
\parallelenpaar{48}{}{132^\circ}{\alpha}{}
\end{minipage}}
\end{aufgabe}

% 9: Winkel im Viereck berechnen – winkel-dreiecke-basis-k4-v2 (Original 2026-FOR-B1i)
\begin{aufgabe}{\begin{minipage}[t]{0.6\linewidth}Ein Parallelogramm hat links unten einen Winkel von $57^\circ$. Wie groß ist der Winkel $\delta$ links oben? \hfill (P26F) \\ $\delta =$ \leerfeld °\end{minipage}\hfill\begin{minipage}[t]{0.37\linewidth}\vspace{0pt}
\parallelogramm[winkel={57^\circ,,,\delta}]{3}{1.8}{57}
\end{minipage}}
\end{aufgabe}

% 10: Wahrscheinlichkeit einstufig – bruchrechnung-basis-k1-v3 (Original 2014-OS-B1d)
\begin{aufgabe}{Ein Glücksrad hat 8 gleich große Felder mit den Zahlen 1 bis 8. Wie groß ist die Wahrscheinlichkeit für eine gerade Zahl? \hfill (P14) \\ \leerfeld}
\end{aufgabe}

\begleitteil
\erg{1}{R, denn $\frac{45}{360} = \frac{1}{8}$}
\erg{2}{$88\,\% = 0{,}88$; die anderen sind $0{,}81$, $0{,}85$ und $0{,}875$}
\erg{3}{$-2 \cdot (-1 - 3) = -2 \cdot (-4) = 8$}
\erg{4}{immer negativ}
\erg{5}{$0{,}4 \cdot 35 = 14\,€$}
\erg{6}{$4 \cdot 1{,}25 = 5\,€$ (nicht nur $1\,€$)}
\erg{7}{$6b + 9$}
\erg{8}{$\alpha = 180^\circ - 132^\circ = 48^\circ$ (Stufenwinkel, dann Nebenwinkel)}
\erg{9}{$\delta = 180^\circ - 57^\circ = 123^\circ$}
\erg{10}{$\frac{4}{8} = \frac{1}{2}$ (2, 4, 6, 8)}
\end{document}
